% Rapport final — Projet DevOps : Flashcards API
% Compilation : pdflatex rapport.tex (deux fois pour la table des matières éventuelle)
\documentclass[11pt,a4paper]{article}

% ---------- Informations de la page de garde (à compléter) ----------
\newcommand{\etablissement}{Nom de l'établissement}
\newcommand{\departement}{Département / Filière}
\newcommand{\module}{DevOps}
\newcommand{\etudiant}{Mouhib Fahem}
\newcommand{\enseignant}{Nom de l'enseignant(e)}
\newcommand{\annee}{2026 -- 2027}
\newcommand{\daterendu}{Octobre 2026}
% --------------------------------------------------------------------

\usepackage[utf8]{inputenc}
\usepackage[T1]{fontenc}
\usepackage[french]{babel}
\usepackage{lmodern}
\usepackage{microtype}
\usepackage[a4paper,margin=1.7cm,top=1.4cm,bottom=1.5cm]{geometry}
\usepackage{xcolor}
\usepackage{tabularx}
\usepackage{booktabs}
\usepackage{enumitem}
\usepackage{tikz}
\usetikzlibrary{arrows.meta,positioning}
\usepackage{titlesec}
\usepackage{fancyhdr}
\usepackage[hidelinks]{hyperref}

\definecolor{accent}{HTML}{0B6BCB}
\definecolor{muted}{HTML}{57606A}
\definecolor{tint}{HTML}{EEF5FC}

\titleformat{\section}{\large\bfseries\color{accent}}{\thesection.}{0.5em}{}[\vspace{-0.6ex}\textcolor{accent!30}{\rule{\linewidth}{0.4pt}}]
\titlespacing*{\section}{0pt}{1.3ex}{0.6ex}
\setlist{nosep,leftmargin=1.3em}
\setlength{\parindent}{0pt}
\setlength{\parskip}{0.45ex}
\renewcommand{\arraystretch}{1.18}

\pagestyle{fancy}
\fancyhf{}
\fancyhead[L]{\small\color{muted}Projet DevOps --- Flashcards API}
\fancyhead[R]{\small\color{muted}\etudiant}
\fancyfoot[C]{\small\color{muted}\thepage}
\renewcommand{\headrulewidth}{0.3pt}

\newcommand{\code}[1]{\texttt{\small #1}}

\begin{document}

% =========================== PAGE DE GARDE ===========================
\begin{titlepage}
  \centering
  {\large\scshape \etablissement\par}
  \vspace{0.2cm}
  {\scshape \departement\par}
  \vspace{0.3cm}
  \textcolor{accent}{\rule{0.85\linewidth}{0.8pt}}\par
  \vspace{3.2cm}

  {\normalsize\scshape\color{muted} Module : \module\par}
  \vspace{0.6cm}
  {\Huge\bfseries Flashcards API\par}
  \vspace{0.4cm}
  {\Large Une petite API, toute la chaîne DevOps\par}
  \vspace{0.3cm}
  {\large\color{muted} CI/CD $\cdot$ Docker $\cdot$ Kubernetes $\cdot$ Observabilité $\cdot$ Sécurité\par}
  \vspace{1.2cm}
  {\Large\bfseries Rapport final de projet\par}

  \vspace{2.6cm}
  \begin{tikzpicture}[node distance=0.35cm, box/.style={draw=accent, fill=tint, rounded corners=3pt, minimum width=2.25cm, minimum height=0.75cm, font=\small}]
    \node[box] (a) {Code};
    \node[box, right=of a] (b) {CI/CD};
    \node[box, right=of b] (c) {Docker};
    \node[box, right=of c] (d) {Kubernetes};
    \node[box, right=of d] (e) {Observabilité};
    \foreach \x/\y in {a/b,b/c,c/d,d/e} \draw[-{Stealth}, muted] (\x) -- (\y);
  \end{tikzpicture}

  \vfill
  \begin{tabular}{@{}ll@{}}
    \textbf{Réalisé par :} & \etudiant \\[0.15cm]
    \textbf{Encadré par :} & \enseignant \\[0.15cm]
    \textbf{Dépôt :} & \url{https://github.com/mouhibfahem/devops} \\[0.15cm]
    \textbf{Image Docker :} & \url{https://hub.docker.com/r/mouhibfahem/flashcards-api} \\
  \end{tabular}

  \vspace{1.2cm}
  \textcolor{accent}{\rule{0.85\linewidth}{0.8pt}}\par
  \vspace{0.2cm}
  {\large Année universitaire \annee \quad$\cdot$\quad \daterendu\par}
\end{titlepage}
% =====================================================================

\section{Objectif et application}
L'objectif était d'appliquer toute la chaîne DevOps à un petit service : du code jusqu'au déploiement, avec
automatisation, observabilité et sécurité. L'application est une \textbf{API REST de flashcards} écrite en Go
(bibliothèque standard, \textbf{149 lignes} dans \code{main.go}) qui aide à réviser des notions du cours.
Elle permet de créer, lister, consulter et supprimer des cartes, et propose un \textbf{mode quiz}
(\code{GET /cards/random}) qui renvoie une question au hasard sans sa réponse. Les tests unitaires couvrent
\textbf{90,8\,\%} du code.

\section{Architecture}
\begin{center}
\resizebox{\linewidth}{!}{%
\begin{tikzpicture}[
    font=\small,
    stage/.style={draw=muted, rounded corners=3pt, font=\footnotesize, minimum width=2.75cm, minimum height=0.8cm, fill=white},
    svc/.style={draw=muted, rounded corners=3pt, minimum width=3.2cm, minimum height=0.85cm, fill=white},
    arr/.style={-{Stealth}, muted},
    lbl/.style={font=\scriptsize, text=muted}]
  % Bande 1 : pipeline
  \fill[gray!6, draw=gray!35, rounded corners=5pt] (-0.25,0.55) rectangle (17.3,-1.15);
  \node[anchor=west, font=\small\bfseries, text=muted] at (-0.1,0.25) {1 $\cdot$ Pipeline GitHub Actions (à chaque Pull Request et merge)};
  \node[stage] (p1) at (1.2,-0.5) {Lint + tests};
  \node[stage] (p2) at (4.7,-0.5) {SAST};
  \node[stage] (p3) at (8.2,-0.5) {Image + Trivy};
  \node[stage] (p4) at (11.7,-0.5) {Kubernetes + DAST};
  \node[stage, draw=green!50!black] (p5) at (15.4,-0.5) {Docker Hub};
  \foreach \x/\y in {p1/p2,p2/p3,p3/p4,p4/p5} \draw[arr] (\x) -- (\y);
  % Bande 2 : exécution
  \fill[gray!6, draw=gray!35, rounded corners=5pt] (-0.25,-1.5) rectangle (17.3,-5.6);
  \node[anchor=west, font=\small\bfseries, text=muted] at (-0.1,-1.8) {2 $\cdot$ Exécution : Kubernetes (kind) ou Docker Compose};
  \node[draw=accent, fill=tint, line width=0.9pt, rounded corners=4pt, minimum width=4.2cm, minimum height=2.9cm,
        align=center] (api) at (2.2,-3.75)
        {\textbf{Flashcards API (Go)}\\[2pt]\scriptsize 2 pods $\cdot$ Service $\cdot$ sondes\\\scriptsize\code{/cards} $\cdot$ \code{/health} $\cdot$ \code{/metrics}\\\scriptsize logs JSON (slog)\\\scriptsize spans OpenTelemetry};
  \node[svc] (prom) at (8.2,-2.85) {Prometheus};
  \node[svc] (graf) at (13.6,-2.85) {Grafana (dashboard)};
  \node[svc] (jae)  at (8.2,-4.55) {Jaeger (traces)};
  \node[svc] (zap)  at (13.6,-4.55) {OWASP ZAP (DAST)};
  \draw[arr] (prom.west) -- node[above, lbl] {scrape} (prom.west -| api.east);
  \draw[arr] (prom) -- (graf);
  \draw[arr] (jae.west -| api.east) -- node[above, lbl] {OTLP} (jae.west);
  \draw[arr] (zap.south) -- ++(0,-0.55) -| node[pos=0.25, below, lbl] {attaque l'API déployée} (api.south);
\end{tikzpicture}%
}
\end{center}
Chaque tâche suit le même flux : une \textbf{issue} (10 issues, tableau Kanban public GitHub Projects), une
\textbf{branche}, une \textbf{Pull Request} liée à l'issue, la vérification par le pipeline, puis le merge.
Au total, \textbf{8 Pull Requests} ont été fusionnées.

\section{Pipeline CI/CD (GitHub Actions)}
\begin{tabularx}{\linewidth}{@{}p{3.3cm}X@{}}
\toprule
\textbf{Job} & \textbf{Ce qu'il fait} \\
\midrule
Lint, test, build & \code{gofmt}, \code{go vet}, tests avec détecteur de concurrence et couverture, compilation \\
SAST & gosec (résultats dans GitHub code scanning), govulncheck (dépendances), gitleaks (secrets dans l'historique) \\
Image & Build Docker unique puis scan \textbf{Trivy} : échec si une faille HIGH/CRITICAL corrigible est présente \\
Deploy + DAST & Cluster \textbf{kind} neuf, déploiement de l'image scannée, tests à travers le Service, scan \textbf{OWASP ZAP} de l'API déployée \\
Publish (main) & Push sur Docker Hub (\code{latest}, \code{sha-<commit>}) \textbf{uniquement si tout est vert} \\
\bottomrule
\end{tabularx}

Les identifiants Docker Hub sont des \textbf{secrets GitHub}, jamais dans le code. Les rapports (couverture,
SAST, Trivy, ZAP) sont conservés comme artefacts de chaque exécution.

\section{Conteneurisation}
\code{Dockerfile} en deux étapes : \code{golang:1.27.2-alpine} compile un binaire statique, puis seul ce binaire
est copié dans une image \textbf{distroless} (environ 30\,Mo, sans shell, utilisateur non-root 65532). Le fichier
\code{docker-compose.yml} lance l'API avec Prometheus, Grafana et Jaeger en une commande.

\section{Observabilité}
\begin{minipage}[t]{0.58\linewidth}
\begin{itemize}
  \item \textbf{Logs} : une ligne JSON par requête (\code{log/slog}) avec méthode, route, statut, durée,
        \code{request\_id}, \code{trace\_id}. Le niveau suit le statut : INFO, WARN (4xx), ERROR (5xx et
        \emph{panics} récupérées).
  \item \textbf{Métriques} : \code{/metrics} au format Prometheus : compteur de requêtes, histogramme des temps de
        réponse, requêtes en cours. Le label \code{route} utilise le motif (\code{GET /cards/\{id\}}) pour limiter
        le nombre de séries.
  \item \textbf{Dashboard} Grafana provisionné automatiquement : requêtes/s, taux d'erreurs, latence p50/p95/p99.
  \item \textbf{Tracing} : un span OpenTelemetry par requête, exporté en OTLP vers Jaeger ; un en-tête
        \code{traceparent} entrant est poursuivi.
\end{itemize}
\end{minipage}\hfill
\begin{minipage}[t]{0.38\linewidth}
\vspace{0pt}
\colorbox{tint}{\parbox{0.94\linewidth}{\small
\textbf{Lien avec le cours SRE.} La latence p95 et le taux d'erreurs sont des \textbf{SLI}. Un \textbf{SLO}
possible : \og 95\,\% des requêtes en moins de 200\,ms et moins de 0,1\,\% d'erreurs 5xx sur un mois \fg. Lors
des tests de charge, la p95 mesurée était inférieure à 1\,ms.}}
\end{minipage}

\section{Sécurité}
\begin{tabularx}{\linewidth}{@{}p{6.2cm}p{2.2cm}X@{}}
\toprule
\textbf{Problème trouvé} & \textbf{Outil} & \textbf{Correction} \\
\midrule
11 vulnérabilités connues dans Go 1.27.0 et \code{x/net} v0.59.0 & govulncheck & Toolchain go1.27.2 et \code{x/net} v0.60.0 \\
Injection dans les logs (G706) & gosec & Attribut structuré \code{slog} \\
En-têtes \code{X-Content-Type-Options} et \code{Cross-Origin-Resource-Policy} absents & OWASP ZAP & En-têtes de sécurité sur toutes les réponses ; le second scan ne garde que des alertes informatives \\
\bottomrule
\end{tabularx}

Autres mesures : aucune fuite de secret (gitleaks), 0 faille HIGH/CRITICAL dans l'image (Trivy), corps de requête
limité à 1\,Mo, délais d'expiration serveur, conteneur non-root et en lecture seule.

\section{Kubernetes}
Les manifests \code{k8s/}, appliqués avec Kustomize, créent un namespace, une ConfigMap, un \textbf{Deployment de
2 réplicas} et un \textbf{Service} ClusterIP. Le Deployment définit des sondes \emph{readiness} et \emph{liveness}
sur \code{/health}, des requêtes et limites CPU/mémoire, une mise à jour progressive sans interruption
(\code{maxUnavailable: 0}) et un contexte de sécurité strict (non-root, système de fichiers en lecture seule,
aucune capability Linux). Test en local avec kind : quand on supprime un pod, Kubernetes le recrée en quelques
secondes.

\section{Leçons apprises}
\begin{itemize}
  \item \textbf{Les scans trouvent de vrais problèmes.} govulncheck a révélé des failles dans la version même de
        Go ; sans automatisation, elles seraient passées inaperçues.
  \item \textbf{L'outillage évolue plus vite que les outils.} gosec ne lisait pas encore Go 1.27 et a dû être
        exécuté avec Go 1.26. Il faut vérifier qu'un outil analyse réellement tout le code.
  \item \textbf{Kubernetes est strict.} \code{runAsNonRoot} exige un UID numérique dans l'image ; l'erreur
        \code{Create\allowbreak Container\allowbreak Config\allowbreak Error} l'a révélé.
  \item \textbf{L'état en mémoire empêche de passer à l'échelle.} Avec 2 réplicas, chaque pod a ses propres
        cartes ; une base de données serait l'étape suivante.
  \item \textbf{Petits incréments.} Une issue = une PR vérifiée par le pipeline : chaque erreur est restée petite et
        facile à corriger.
\end{itemize}

\section{Pistes d'amélioration}
Base de données (PostgreSQL) pour partager les cartes entre pods, déploiement sur un cloud (bonus), Prometheus et
Grafana dans le cluster, alertes basées sur les SLO, et spans enfants pour un tracing plus détaillé.

\end{document}
